%!TEX root = ../main.tex
\chapter{Conclusions}
\label{ch:conclusion}

We presented Adaptive-$K$, a unified frequency-domain defense for
protecting deep-learning-based automatic modulation classifiers against
adversarial attacks. The method combines spectral-flatness routing,
SNR-adaptive K-cap, and per-sample magnitude-knee estimation in a single
shared-FFT pass, requires no classifier retraining, and adds
sub-millisecond overhead per batch of 64 samples in our GPU implementation,
making it compatible with
low-latency receiver-side preprocessing for spectrum monitoring and
cognitive radio applications.

Through systematic evaluation on RML2016.10a against five attacks (CW,
EAD-L1, EAD-EN, FGSM, PGD) and comparison with five per-SNR-calibrated
classical signal-processing baselines and an empirical
randomized-smoothing baseline ($k = 20$ noisy copies at $\sigma = 0.01$;
no formal certificate), we established four principal findings:

\begin{enumerate}
\item \textbf{Adversarial attacks on AMC behave as control-plane attacks.}
  A two-track CRC analysis on the strongest optimization-based attack
  (CW) shows that adversarial perturbations collapse AMC accuracy to
  near zero while leaving the underlying signal data largely decodable
  (CRC pass rate $> 71\%$ with oracle demodulation) at SNR 18\,dB
  across eight digital modulations. This characterization motivates
  recovery as the correct defense strategy: the perturbation need only be
  attenuated for correct classification to be restored. Extending this
  CRC analysis to EAD, FGSM, and PGD is future work.

\item \textbf{Adaptive-$K$ outperforms all classical baselines across
  attack families.} The unified method achieves the highest recovery
  accuracy on four of five attacks (all gains reported over undefended
  and over the best classical baseline on that attack): 71.3\% on CW
  (+35.9\,pp / +6.2\,pp over Kalman), 70.8\% on EAD-L1
  (+64.2\,pp / +8.5\,pp over Kalman), 67.5\% on EAD-EN
  (+64.1\,pp / +7.3\,pp over Kalman), and 43.4\% on PGD
  (+10.9\,pp / +4.6\,pp over FIR). The SNR-adaptive K-cap ensures
  robust performance across the full SNR range by automatically adjusting
  filtering aggressiveness to estimated signal quality.

\item \textbf{Classical signal-processing filters are counterproductive
  on gradient attacks.} On FGSM, all five classical filters reduce
  accuracy below the undefended baseline (48.1\%) despite per-SNR
  calibration. This negative result indicates that deploying conventional
  filtering as a security measure can provide a false sense of protection:
  gradient-sign perturbations are spectrally correlated with the signal,
  so time-domain smoothing removes signal energy alongside perturbation
  energy.

\item \textbf{Wideband routing is essential for universal defense.}
  Without spectral-flatness routing, Top-$K$ filtering destroys AM-SSB
  signals. The flatness-based routing in Adaptive-$K$ correctly
  identifies wideband samples and applies per-sample quantization instead,
  recovering AM-SSB accuracy to 94\%.
\end{enumerate}

Several directions remain open for future work.

\textbf{Adaptive defense-aware attacks.}
The evaluation in this thesis assumes a defense-unaware attacker. A
white-box adversary with knowledge of the Top-$K$ operator could craft
perturbations that survive the frequency mask. Evaluating such
\emph{adaptive defense-aware attacks}---in which the attack optimization
incorporates the defense---is an important robustness validation that we
explicitly do not claim coverage of in this work, and is the most
important next step.

\textbf{Generalization to longer and richer signals.}
Extending the evaluation to RML2018.01a (1024-sample signals, 24 classes)
would establish the generality of the proposed defense across longer
signals and a richer modulation set, and would test whether the
SNR-adaptive K-cap continues to hold without retuning at the longer
$T = 1024$ FFT length.

\textbf{Combination with adversarial training.}
Combining Adaptive-$K$ with adversarial training could provide
complementary robustness: training-time regularization hardens the
classifier's decision boundaries while inference-time spectral filtering
removes perturbation energy before the signal reaches the network. A
particularly interesting question is whether such a hybrid pipeline
remains effective against adaptive defense-aware attacks.

\textbf{Certified randomized smoothing.}
The randomized-smoothing baseline used in this thesis is empirical:
it runs the Cohen et al.\ majority-vote rule over $k = 20$ noisy
copies at $\sigma = 0.01$ (via either the in-tree implementation or
the IBM ART~\cite{art2018} wrapper, which are numerically
equivalent), and does not compute the Clopper--Pearson lower bound on
$p_A$ required for a certified $L_2$ radius. Evaluating
a fully certified configuration (hundreds to thousands of noisy copies
per sample, with formal certificates) would clarify the trade-off
between certified and empirical defenses for AMC and quantify the
throughput cost of certification relative to Adaptive-$K$.
